% ============================================================
%  Block B -- on an interval, and integrability
% ============================================================
\rsection{Block B --- on an interval}

\begin{signpost}[What changes when the point is dropped]
Promote each local property to ``at every point of $[a,b]$'' and a fourth
property joins them: Riemann integrability, which has no pointwise form at
all. The arrow into it does not come from anything local --- it comes from
compactness.
\end{signpost}

\ritem{T7}{Theorem}
\emph{If $f$ is differentiable at every point of $[a,b]$ then $f$ is
continuous on $[a,b]$.}

\begin{proof}
T1 at each point.
\end{proof}

\ritem{T8}{Theorem}
\emph{If $f$ is continuous on $[a,b]$ then $f$ is Riemann integrable there.}

\begin{proof}
Let $\varepsilon>0$. Since $[a,b]$ is compact, $f$ is \emph{uniformly}
continuous (4.19), so there is $\delta>0$ with
$|f(x)-f(y)| < \varepsilon/(b-a)$ whenever $|x-y|<\delta$. Take any partition
$P$ of mesh less than $\delta$. On each subinterval $[x_{i-1},x_i]$ the
supremum $M_i$ and infimum $m_i$ are attained (4.16), at points less than
$\delta$ apart, so $M_i - m_i \le \varepsilon/(b-a)$ and
\[ U(P,f) - L(P,f) = \sum_i (M_i-m_i)\,\Delta x_i
   \;\le\; \frac{\varepsilon}{b-a}\sum_i \Delta x_i = \varepsilon . \]
By Rudin's criterion 6.6, $f$ is integrable.
\end{proof}

\begin{proofwalk}[How the proof flows]
  \item \textbf{Choose the method.} Integrability is a statement about upper
        and lower sums being close, so the target is $U-L<\varepsilon$. That
        is a statement about the \emph{oscillation} of $f$ on the pieces of a
        partition, so the whole problem is: make the oscillation uniformly
        small.
  \item \textbf{Spend compactness --- once.} Plain continuity gives a
        $\delta$ at each point, and those $\delta$'s could shrink to nothing
        as the point moves. Uniform continuity is exactly the statement that
        one $\delta$ serves everywhere, and 4.19 is where compactness is
        spent. Everything after this is bookkeeping.
  \item \textbf{Convert to the partition.} Mesh smaller than $\delta$ makes
        every subinterval a place where the uniform estimate applies.
  \item \textbf{Sum.} The $\Delta x_i$ add to $b-a$, which is why the
        $\varepsilon/(b-a)$ was chosen in step 2 --- reverse-engineered from
        the final line, the same way Rudin picks his constants at 1.21.
  \item \textbf{Audit.} Compactness is spent once, in step 2; the
        attainment of $M_i$ and $m_i$ (4.16) is a convenience, not a
        necessity --- $\sup - \inf \le \varepsilon/(b-a)$ would follow from
        the uniform estimate anyway. Note what is \emph{not} used: nothing
        about differentiability, and nothing pointwise beyond continuity.
\end{proofwalk}

\rsection{Which functions are integrable}

\begin{signpost}[From a gallery to a theorem]
The examples below could each be settled by hand with upper and lower sums.
Doing so leaves the impression that integrability is a matter of luck. It is
not: there is an exact criterion, and it is stated in the language Chapter
2's and Chapter 4's appendices already built.
\end{signpost}

\ritem{T9}{Definition}
\emph{A set $E \subset \R$ has \textbf{measure zero} if for every
$\varepsilon>0$ there are countably many intervals $I_1,I_2,\dots$ with
$E \subset \bigcup_k I_k$ and $\sum_k |I_k| < \varepsilon$.}

\begin{center}
\begin{tikzpicture}[x=1mm, y=1mm]
  \draw[axis] (0,0) -- (104,0);
  \foreach \x in {8,14,17,31,44,46,60,73,74,88,95}{
    \node[ptfill, minimum size=2.6pt] at (\x,0) {};}
  \foreach \a/\b in {6.4/9.6, 13/18.2, 29.8/32.2, 43/47.2, 59/61, 72/75.4,
                     87.2/88.8, 94.4/95.6}{
    \draw[thin] (\a,4) -- (\b,4);
    \draw[thin] (\a,3.2) -- (\a,4.8);
    \draw[thin] (\b,3.2) -- (\b,4.8);
    \draw[guide] (\a,0.6) -- (\a,3.2);
    \draw[guide] (\b,0.6) -- (\b,3.2);}
  \node[lbl, anchor=west] at (0,10)
    {countably many intervals, total length $<\varepsilon$};
  \node[lbl, anchor=west] at (0,-6)
    {the set $E$ --- may be infinite, even dense};
\end{tikzpicture}
\end{center}
\figcap{Measure zero. The intervals may be as numerous as you like, so long
as their lengths add to less than $\varepsilon$ --- and $\varepsilon$ is
arbitrary. Every countable set qualifies: cover its $k$th point by an
interval of length $\varepsilon/2^{k+1}$. No interval $[a,b]$ with $a<b$
qualifies.}

\ritem{T10}{Theorem}
\emph{(Lebesgue's criterion.) A bounded $f$ on $[a,b]$ is Riemann integrable
if and only if its set of discontinuities has measure zero.}

\begin{proof}[Proof sketch]
Recall from S2.5 and S4.5 that with
$\Omega_\varepsilon = \{x : \omega_f(x) \ge \varepsilon\}$, each
$\Omega_\varepsilon$ is \emph{closed} --- hence compact here --- and the
discontinuity set is $D = \bigcup_n \Omega_{1/n}$.

($\Rightarrow$) Let $f$ be integrable and fix $n$. Given $\eta>0$ choose a
partition with $U-L < \eta/n$. Every subinterval whose interior meets
$\Omega_{1/n}$ contributes at least $(1/n)\Delta x_i$ to $U-L$, so those
subintervals have total length less than $\eta$. They cover
$\Omega_{1/n}$ apart from finitely many partition points, which are covered
by intervals of arbitrarily small total length. So each $\Omega_{1/n}$ has
measure zero, and $D$, a countable union of such sets, does too.

($\Leftarrow$) Suppose $D$ has measure zero and $|f| \le M$. Fix
$\varepsilon>0$. Then $\Omega_{\varepsilon}$ is compact and of measure zero,
so finitely many open intervals of total length $<\varepsilon$ cover it. Off
their union, every point has $\omega_f < \varepsilon$, hence a neighbourhood
on which the oscillation is below $\varepsilon$; compactness of the remaining
closed set gives a finite subcover and so a partition on which those pieces
each oscillate by less than $\varepsilon$. Splitting $U-L$ into the two kinds
of subinterval bounds it by $2M\varepsilon + \varepsilon(b-a)$, which is
arbitrarily small.
\end{proof}

\begin{deeper}[The criterion settles the whole row at once]
Everything in the integrability row now follows by inspection of a
discontinuity set:

\smallskip
\begin{tabular}{@{}l l l@{}}
\toprule
\textbf{function} & \textbf{discontinuous on} & \textbf{integrable?} \\
\midrule
step function (E5)   & a finite set        & yes \\
Thomae (E6)          & $\Q$ --- countable  & yes \\
Dirichlet (E7)       & everything          & no \\
$x^2\sin(1/x)$ (E3)  & nowhere             & yes \\
\bottomrule
\end{tabular}

\smallskip
Two things are worth dwelling on. First, E6 is integrable while being
discontinuous on a \emph{dense} set --- so ``integrable'' is compatible with
very bad pointwise behaviour, and there is no arrow at all from
integrability back to continuity. Second, the machinery this proof runs on
is not new: the closedness of $\Omega_\varepsilon$ and the decomposition
$D = \bigcup_n \Omega_{1/n}$ are S2.5 and S4.5, proved in the appendices for
quite different reasons. Oscillation was the right idea twice.
\end{deeper}

\ritem{T11}{Theorem}
\emph{Bounded does not imply integrable, and integrable does not imply
continuous anywhere in particular.}

\begin{proof}
E7 is bounded and discontinuous everywhere, so by T10 it is not integrable.
E6 is integrable by T10 and is discontinuous at every rational.
\end{proof}

\rsection{The deepest cell: a derivative that is not integrable}

\ritem{T12}{Theorem}
\emph{(Volterra.) There is a function $F$ on $[0,1]$, differentiable at every
point, whose derivative $F'$ is bounded but \textbf{not} Riemann integrable.}

\begin{deeper}[The construction, and what it costs the Fundamental Theorem]
The construction starts from a \emph{fat Cantor set} $K$ --- built like
Rudin's 2.44 but removing middle intervals of rapidly shrinking length, so
that $K$ is nowhere dense yet does \emph{not} have measure zero. On each
interval removed one plants a scaled copy of $x^2\sin(1/x)$ (E3), tapered so
that it and its derivative vanish at the endpoints, and sets $F=0$ on $K$.

The resulting $F$ is differentiable everywhere, with $|F'| \le 1$. But at
every point of $K$ the derivative oscillates: arbitrarily close by there are
points where $F'$ takes values near $\pm1$. So $\Omega_1 \supset K$, and $K$
does not have measure zero --- by T10, $F'$ is not integrable.

\smallskip
\textbf{What this costs.} Rudin's 6.21 says: if $f$ is integrable on $[a,b]$
and $F'=f$, then $\int_a^b f = F(b)-F(a)$. Volterra's example shows the
hypothesis ``$f$ is integrable'' cannot be dropped --- there really are
derivatives to which the formula does not apply, because the integral on the
left does not exist. So the two halves of the Fundamental Theorem are
independent statements: having an antiderivative does not make a function
integrable, and being integrable does not make a function a derivative
(E5, a step function, is integrable and is not a derivative, by T4).
\end{deeper}

\begin{pitfall}[Four properties that sound alike]
Keep these apart. For a bounded $f$ on $[a,b]$:
\begin{enumerate}[label=(\alph*), leftmargin=2.2em, itemsep=2pt, topsep=3pt]
  \item $f$ is continuous;
  \item $f$ is Riemann integrable;
  \item $f$ has an antiderivative;
  \item $f$ is a derivative \emph{and} integrable, so that 6.21 applies.
\end{enumerate}
(a) implies all the others. (b) and (c) are independent: E5 is integrable
with no antiderivative (T4), and Volterra's $F'$ has an antiderivative and is
not integrable. (d) is their conjunction, and it is what the Fundamental
Theorem actually requires.
\end{pitfall}
